\documentclass[11pt,letterpaper]{article}
\usepackage[margin=0.7in,headheight=15pt]{geometry}
\usepackage[T1]{fontenc}\usepackage{lmodern,amsmath,amssymb,fancyhdr,parskip,hyperref}
\hypersetup{colorlinks=true,urlcolor=blue}
\pdfinfoomitdate=1\pdftrailerid{}\pdfsuppressptexinfo=15
\pagestyle{fancy}\fancyhf{}\lhead{Week 82 / Endless class lines}\rhead{Facilitator draft}\cfoot{Bellingham Math Circle / Week 82 / facilitator / \thepage}
\setlength{\parskip}{7pt}\setlength{\parindent}{0pt}
\begin{document}
\section*{Mathematical overview}
A complete matching (bijection) uses every kid exactly once. An order-preserving matching additionally preserves before/after relations. Labels identify kids, not required partners; equal labels in two class lines impose no matching constraint. For the endless class line $1,2,3,\ldots$, prepending a star admits the matching star$\mapsto1$, $n\mapsto n+1$. Thus the order type $1+\omega$ is $\omega$. Appending a star after \emph{every} numbered kid produces a last element; the original class line has none. Hence $\omega+1$ is a different order type.

Two consecutive endless blocks have type $\omega+\omega$. The first kid in the second block has infinitely many predecessors and no immediate predecessor. An endless row of pairs has type $2\cdot\omega=\omega$; every kid has finitely many predecessors. They cannot be matched preserving order. All these class lines have the same countable cardinal size, using explicit matches that need not preserve order.

Children should discover a structural obstruction or a rule covering every kid. Finite drawn matches are experiments; a uniform assignment and inverse establish an infinite matching. Grades 3--5 only, with facility interpreting a finite schematic for an endless rule. Ordinal addition here means concatenation. Surreal addition is a distinct commutative operation, so do not transfer these addition identities to Week 83's numbers.
\newpage
\section*{Preparation, launch and readiness}
Per pair: numbered kid cards (student page 5), stars and after-every-kid panels (page 6), blank paper for class lines and yarn or dry-erase arrows. For endless rows, draw a few numbered kids followed by an arrow marked ``numbers continue forever,'' as on student pages 2--3; use two colors when needed. A continuation symbol is an instruction, never another kid. A star after every number uses a \emph{separate panel} so children do not read it as after just the printed prefix. Colors and labels identify kids; they do not restrict which kids may be partners. Problem 4's small boxes are sketch space: arrange the cards and hand-drawn endless rows on the table, then sketch the result in those boxes. Provide sufficient table space for two boards together; enlarging to two joined Letter sheets is possible. Physical fit untested.

Prerequisites: compare positions, follow a matching rule, understand ``each kid has exactly one partner.'' Reading is brief and can be read aloud. Reasoning load lies in extending a rule to an arbitrary named position, not counting a printed finite list. One anchored adult can coach pairs; at the upper triplet rotate constructor/checker/questioner.

Launch together: let children reorder a few finite cards, then compare the three printed examples: five-to-five in order; five-to-five with the first two partners exchanged; and five-to-six, where one kid must be left out. The last picture shows one failed attempt, and counting proves every attempt fails. Ask the partner to check both unused kids and crossing order. Introduce the continuation-arrow convention \emph{before} asking about an endless class line.

First route after movement: ten minutes on finite Problem 1, 15 minutes prepending and matching, 15 minutes appending and explaining its obstruction. Problem 3's two blocks versus pairs is a second visit or encore; Problem 4 is open construction when the basic convention is controlled. Ordinal names may be introduced after the need arises. Do not require all stations or notation in one hour. Whole-group launch, then selected participating tables only.
\newpage
\section*{Facilitator explanations and hints}
Problem 1: list the B partners of R1, R2, R3. The six lists are\[ (1,2,3),(1,3,2),(2,1,3),(2,3,1),(3,1,2),(3,2,1). \] Only the first keeps every pair in order. Exhaustiveness follows from three choices, then two, then one.

Problem 2, prepending: the assignment is onto because kid 1 receives the star and kid $m\ge2$ receives original $m-1$. It is one-to-one and preserves order. Matching only six physical cards leaves an apparent unmatched card at the right edge; the repeating rule, not this artificial drawing boundary, resolves the infinite case.

Appending: suppose there were an order-preserving bijection to the ordinary endless class line. The image of the appended last kid would have to be last, but every ordinary numbered kid has a successor. Contradiction. The issue is not inability to physically walk past infinitely many kids.

Ignoring order, star$\mapsto1$, $n\mapsto n+1$ still gives a size match for the appended-star class line. It moves the last kid to the front, showing why the preservation condition matters. It does not make the class lines have equal order type.

Problem 3: D is drawn on one horizontal line, with all R kids followed by all B kids. The gap separates entire endless blocks, not just the three printed red cards; E interleaves R and B. Matching each named R$_n$ and B$_n$ in D to its namesake in E gives an unrestricted size match. Their positions in E are $2n-1$ and $2n$. This interleaves blocks and violates their original order. No order-preserving match can send B$_1$ to a kid with only finitely many predecessors, since B$_1$ has all R kids before it.

Problem 4 examples: a row alone versus a star followed by the row admits the order-preserving shift of Problem 2. A row alone versus a row followed by a star admits a size match but no order match because of the last kid. Each construction uses one row and at most two stars.

Hints: point to a last kid if one exists; ask about the first kid of block B; ask how many kids are before a \emph{named} position. Avoid directing every match in advance. The representation is conceptual; drawn endless rows need not physically contain every kid.
\newpage
\section*{Encore, sources and limits}
Encore: attach two, then three kids before or after the endless class line. Compare removing the first kid with removing an appended last kid. Later arrange endlessly many endless blocks as rows with explicit row order; this opens a doorway toward $\omega^2$ but is not required for the present proof.

Sources: Joel David Hamkins, ``Counting to Infinity and Beyond'' (September 2, 2021), \url{https://jdh.hamkins.org/counting-to-infinity-poster/}, for ordinal visualization and continuation to $\omega^2$. Our matching boards and explanatory sequence are independently authored; no poster or video frame is republished. Basic order-type arguments are given fully in this guide.

Pedagogy: \emph{Math Circle by the Bay}, preface p. viii, connects accessible topics to deep mathematics; pp. ix--x discuss materials and varying depth. The grade gate here is our unpiloted design judgment. Existing Week 68 endless graphs explores ends, not these class line order types.

Unpiloted. Print US Letter at Actual Size. Physical board fit, yarn handling, table-space requirement, timing and classroom procedure have not been rehearsed. Record whether children interpreted continuation as a rule, whether they distinguished size/order, and any useful invented class lines. A decorative infinity symbol alone is not evidence that the child controlled the investigation.
\end{document}
